\documentclass[11pt,a4paper]{article}
\usepackage[T1]{fontenc}
\usepackage[utf8]{inputenc}
\usepackage[english]{babel}
\usepackage{amsmath,amssymb}
\usepackage[margin=2cm]{geometry}
\setlength{\parindent}{0pt}
\setlength{\emergencystretch}{3em}
\allowdisplaybreaks

\title{About \texttt{worked\_solutions.tex}}
\author{}
\date{}
\begin{document}
\maketitle
This note explains what the test data in \texttt{worked\_solutions.tex} is used for. That file is input for the engine in \texttt{main.c}; it is read by the test \texttt{test\_worked\_solutions.c} in this folder.

\begin{enumerate}
\item \textbf{Purpose.} The file holds 12 statements decided by the sign table: requests $\{x \in \mathbb{R} \mid \ldots\}$, claims ``for every $x$'' and ``there is an $x$'', and claims that a condition has given solutions. The verdicts are checked elsewhere. Here the engine must also record how each one is reached, one step at a time, so that the final document can explain it as a teacher would.

\item \textbf{What it exercises.} The explanation layer (chapter 40 of \texttt{main.c}, chapter 41 of \texttt{output.c}): the difference of the two sides, the same written as one fraction of polynomials, numerator and denominator split into the factors of the table, the zeros of each factor, the sign table with one line for each factor, the two sides computed at one value of each interval, and the answer read off the last line of the table. Every recorded step is checked again by reading back the text that will be printed.

\item \textbf{Variety.} A chain with all three steps, $\frac{x - 2}{x + 3} \ge x$; one with only the factors, $x^{2} - 5x + 6 < 0$; one with nothing to rewrite; a factor that is never zero; zeros written with radicals; a point where a side has no value; a table with a single column; a restricted set, $\mathbb{R}^{+}$; two false equivalences, refuted by one column; a power of a factor; a system of two comparisons.

\item \textbf{Expected result.} The line \texttt{PASS: 170 of 170 checks on worked solutions passed}. The expected steps, signs and answers are written in the test, with the computations behind them in a comment; they were worked out by hand, not copied from the engine's output.

\item \textbf{How to run it.} From the top folder of the repository:
\begin{verbatim}
cc -std=c11 -O2 -Wall -Wextra -o /tmp/test_worked_solutions \
    tests/equations_inequalities/test_worked_solutions.c \
    && /tmp/test_worked_solutions
\end{verbatim}
\end{enumerate}
\end{document}
